% ═══════════════════════════════════════════════════════════════════════
\section{Des monstres aux fractales}
% ═══════════════════════════════════════════════════════════════════════

\begin{frame}{Ce que les deux premiers chapitres ont préparé}
  \small
  \begin{itemize}
    \item Chapitre 1 : une transformation affine du plan tient en six
      nombres ; similitudes et contractions.
    \item Chapitre 2 : l'espace $\big(\Comp{\R^2},\dH\big)$ des
      compacts du plan, et le théorème du point fixe de Banach.
    \item Aujourd'hui : quelques transformations affines suffisent à
      définir un flocon au contour infiniment dentelé, ou une fougère.
    \item L'existence de ces figures sera justifiée par Banach,
      appliqué à l'espace des compacts.
  \end{itemize}
\end{frame}

\begin{frame}{Figure 3.1 --- Le flocon de von Koch}
  \begin{columns}[c]
    \begin{column}{0.45\textwidth}
      \FigBox[0.66\textheight]{%
        \begin{tikzpicture}[scale=4.4]
          \fill[steelblue4!14] (0,0) -- (0.5,0.8660254) -- (1,0) -- cycle;
          \begin{scope}[fractale/.style={fill=steelblue4!22, draw=sectioncolor, thin, line join=round}]
            % Flocon de von Koch, étape 4
% Fichier généré par generer.py : ne pas modifier à la main.
\fill[fractale]
  (0,0) -- (0.0062,0.0107) -- (0,0.0214) -- (0.0123,0.0214) -- (0.0185,0.0321) -- (0.0123,0.0428) -- (0,0.0428) -- (0.0062,0.0535) -- (0,0.0642) -- (0.0123,0.0642) -- (0.0185,0.0748) -- (0.0247,0.0642) -- (0.037,0.0642) -- (0.0432,0.0748) -- (0.037,0.0855) -- (0.0494,0.0855) -- (0.0556,0.0962) -- (0.0494,0.1069) -- (0.037,0.1069) -- (0.0432,0.1176) -- (0.037,0.1283) -- (0.0247,0.1283) -- (0.0185,0.1176) -- (0.0123,0.1283) -- (0,0.1283) -- (0.0062,0.139) -- (0,0.1497) -- (0.0123,0.1497) -- (0.0185,0.1604) -- (0.0123,0.1711) -- (0,0.1711) -- (0.0062,0.1818) -- (0,0.1925) -- (0.0123,0.1925) -- (0.0185,0.2031) -- (0.0247,0.1925) -- (0.037,0.1925) -- (0.0432,0.2031) -- (0.037,0.2138) -- (0.0494,0.2138) -- (0.0556,0.2245) -- (0.0617,0.2138) -- (0.0741,0.2138) -- (0.0679,0.2031) -- (0.0741,0.1925) -- (0.0864,0.1925) -- (0.0926,0.2031) -- (0.0988,0.1925) -- (0.1111,0.1925) -- (0.1173,0.2031) -- (0.1111,0.2138) -- (0.1235,0.2138) -- (0.1296,0.2245) -- (0.1235,0.2352) -- (0.1111,0.2352) -- (0.1173,0.2459) -- (0.1111,0.2566) -- (0.1235,0.2566) -- (0.1296,0.2673) -- (0.1358,0.2566) -- (0.1481,0.2566) -- (0.1543,0.2673) -- (0.1481,0.278) -- (0.1605,0.278) -- (0.1667,0.2887) -- (0.1605,0.2994) -- (0.1481,0.2994) -- (0.1543,0.3101) -- (0.1481,0.3208) -- (0.1358,0.3208) -- (0.1296,0.3101) -- (0.1235,0.3208) -- (0.1111,0.3208) -- (0.1173,0.3314) -- (0.1111,0.3421) -- (0.1235,0.3421) -- (0.1296,0.3528) -- (0.1235,0.3635) -- (0.1111,0.3635) -- (0.1173,0.3742) -- (0.1111,0.3849) -- (0.0988,0.3849) -- (0.0926,0.3742) -- (0.0864,0.3849) -- (0.0741,0.3849) -- (0.0679,0.3742) -- (0.0741,0.3635) -- (0.0617,0.3635) -- (0.0556,0.3528) -- (0.0494,0.3635) -- (0.037,0.3635) -- (0.0432,0.3742) -- (0.037,0.3849) -- (0.0247,0.3849) -- (0.0185,0.3742) -- (0.0123,0.3849) -- (0,0.3849) -- (0.0062,0.3956) -- (0,0.4063) -- (0.0123,0.4063) -- (0.0185,0.417) -- (0.0123,0.4277) -- (0,0.4277) -- (0.0062,0.4384) -- (0,0.4491) -- (0.0123,0.4491) -- (0.0185,0.4597) -- (0.0247,0.4491) -- (0.037,0.4491) -- (0.0432,0.4597) -- (0.037,0.4704) -- (0.0494,0.4704) -- (0.0556,0.4811) -- (0.0494,0.4918) -- (0.037,0.4918) -- (0.0432,0.5025) -- (0.037,0.5132) -- (0.0247,0.5132) -- (0.0185,0.5025) -- (0.0123,0.5132) -- (0,0.5132) -- (0.0062,0.5239) -- (0,0.5346) -- (0.0123,0.5346) -- (0.0185,0.5453) -- (0.0123,0.556) -- (0,0.556) -- (0.0062,0.5667) -- (0,0.5774) -- (0.0123,0.5774) -- (0.0185,0.588) -- (0.0247,0.5774) -- (0.037,0.5774) -- (0.0432,0.588) -- (0.037,0.5987) -- (0.0494,0.5987) -- (0.0556,0.6094) -- (0.0617,0.5987) -- (0.0741,0.5987) -- (0.0679,0.588) -- (0.0741,0.5774) -- (0.0864,0.5774) -- (0.0926,0.588) -- (0.0988,0.5774) -- (0.1111,0.5774) -- (0.1173,0.588) -- (0.1111,0.5987) -- (0.1235,0.5987) -- (0.1296,0.6094) -- (0.1235,0.6201) -- (0.1111,0.6201) -- (0.1173,0.6308) -- (0.1111,0.6415) -- (0.1235,0.6415) -- (0.1296,0.6522) -- (0.1358,0.6415) -- (0.1481,0.6415) -- (0.1543,0.6522) -- (0.1481,0.6629) -- (0.1605,0.6629) -- (0.1667,0.6736) -- (0.1728,0.6629) -- (0.1852,0.6629) -- (0.179,0.6522) -- (0.1852,0.6415) -- (0.1975,0.6415) -- (0.2037,0.6522) -- (0.2099,0.6415) -- (0.2222,0.6415) -- (0.216,0.6308) -- (0.2222,0.6201) -- (0.2099,0.6201) -- (0.2037,0.6094) -- (0.2099,0.5987) -- (0.2222,0.5987) -- (0.216,0.588) -- (0.2222,0.5774) -- (0.2346,0.5774) -- (0.2407,0.588) -- (0.2469,0.5774) -- (0.2593,0.5774) -- (0.2654,0.588) -- (0.2593,0.5987) -- (0.2716,0.5987) -- (0.2778,0.6094) -- (0.284,0.5987) -- (0.2963,0.5987) -- (0.2901,0.588) -- (0.2963,0.5774) -- (0.3086,0.5774) -- (0.3148,0.588) -- (0.321,0.5774) -- (0.3333,0.5774) -- (0.3395,0.588) -- (0.3333,0.5987) -- (0.3457,0.5987) -- (0.3519,0.6094) -- (0.3457,0.6201) -- (0.3333,0.6201) -- (0.3395,0.6308) -- (0.3333,0.6415) -- (0.3457,0.6415) -- (0.3519,0.6522) -- (0.358,0.6415) -- (0.3704,0.6415) -- (0.3765,0.6522) -- (0.3704,0.6629) -- (0.3827,0.6629) -- (0.3889,0.6736) -- (0.3827,0.6843) -- (0.3704,0.6843) -- (0.3765,0.695) -- (0.3704,0.7057) -- (0.358,0.7057) -- (0.3519,0.695) -- (0.3457,0.7057) -- (0.3333,0.7057) -- (0.3395,0.7163) -- (0.3333,0.727) -- (0.3457,0.727) -- (0.3519,0.7377) -- (0.3457,0.7484) -- (0.3333,0.7484) -- (0.3395,0.7591) -- (0.3333,0.7698) -- (0.3457,0.7698) -- (0.3519,0.7805) -- (0.358,0.7698) -- (0.3704,0.7698) -- (0.3765,0.7805) -- (0.3704,0.7912) -- (0.3827,0.7912) -- (0.3889,0.8019) -- (0.3951,0.7912) -- (0.4074,0.7912) -- (0.4012,0.7805) -- (0.4074,0.7698) -- (0.4198,0.7698) -- (0.4259,0.7805) -- (0.4321,0.7698) -- (0.4444,0.7698) -- (0.4506,0.7805) -- (0.4444,0.7912) -- (0.4568,0.7912) -- (0.463,0.8019) -- (0.4568,0.8126) -- (0.4444,0.8126) -- (0.4506,0.8233) -- (0.4444,0.834) -- (0.4568,0.834) -- (0.463,0.8446) -- (0.4691,0.834) -- (0.4815,0.834) -- (0.4877,0.8446) -- (0.4815,0.8553) -- (0.4938,0.8553) -- (0.5,0.866) -- (0.5062,0.8553) -- (0.5185,0.8553) -- (0.5123,0.8446) -- (0.5185,0.834) -- (0.5309,0.834) -- (0.537,0.8446) -- (0.5432,0.834) -- (0.5556,0.834) -- (0.5494,0.8233) -- (0.5556,0.8126) -- (0.5432,0.8126) -- (0.537,0.8019) -- (0.5432,0.7912) -- (0.5556,0.7912) -- (0.5494,0.7805) -- (0.5556,0.7698) -- (0.5679,0.7698) -- (0.5741,0.7805) -- (0.5802,0.7698) -- (0.5926,0.7698) -- (0.5988,0.7805) -- (0.5926,0.7912) -- (0.6049,0.7912) -- (0.6111,0.8019) -- (0.6173,0.7912) -- (0.6296,0.7912) -- (0.6235,0.7805) -- (0.6296,0.7698) -- (0.642,0.7698) -- (0.6481,0.7805) -- (0.6543,0.7698) -- (0.6667,0.7698) -- (0.6605,0.7591) -- (0.6667,0.7484) -- (0.6543,0.7484) -- (0.6481,0.7377) -- (0.6543,0.727) -- (0.6667,0.727) -- (0.6605,0.7163) -- (0.6667,0.7057) -- (0.6543,0.7057) -- (0.6481,0.695) -- (0.642,0.7057) -- (0.6296,0.7057) -- (0.6235,0.695) -- (0.6296,0.6843) -- (0.6173,0.6843) -- (0.6111,0.6736) -- (0.6173,0.6629) -- (0.6296,0.6629) -- (0.6235,0.6522) -- (0.6296,0.6415) -- (0.642,0.6415) -- (0.6481,0.6522) -- (0.6543,0.6415) -- (0.6667,0.6415) -- (0.6605,0.6308) -- (0.6667,0.6201) -- (0.6543,0.6201) -- (0.6481,0.6094) -- (0.6543,0.5987) -- (0.6667,0.5987) -- (0.6605,0.588) -- (0.6667,0.5774) -- (0.679,0.5774) -- (0.6852,0.588) -- (0.6914,0.5774) -- (0.7037,0.5774) -- (0.7099,0.588) -- (0.7037,0.5987) -- (0.716,0.5987) -- (0.7222,0.6094) -- (0.7284,0.5987) -- (0.7407,0.5987) -- (0.7346,0.588) -- (0.7407,0.5774) -- (0.7531,0.5774) -- (0.7593,0.588) -- (0.7654,0.5774) -- (0.7778,0.5774) -- (0.784,0.588) -- (0.7778,0.5987) -- (0.7901,0.5987) -- (0.7963,0.6094) -- (0.7901,0.6201) -- (0.7778,0.6201) -- (0.784,0.6308) -- (0.7778,0.6415) -- (0.7901,0.6415) -- (0.7963,0.6522) -- (0.8025,0.6415) -- (0.8148,0.6415) -- (0.821,0.6522) -- (0.8148,0.6629) -- (0.8272,0.6629) -- (0.8333,0.6736) -- (0.8395,0.6629) -- (0.8519,0.6629) -- (0.8457,0.6522) -- (0.8519,0.6415) -- (0.8642,0.6415) -- (0.8704,0.6522) -- (0.8765,0.6415) -- (0.8889,0.6415) -- (0.8827,0.6308) -- (0.8889,0.6201) -- (0.8765,0.6201) -- (0.8704,0.6094) -- (0.8765,0.5987) -- (0.8889,0.5987) -- (0.8827,0.588) -- (0.8889,0.5774) -- (0.9012,0.5774) -- (0.9074,0.588) -- (0.9136,0.5774) -- (0.9259,0.5774) -- (0.9321,0.588) -- (0.9259,0.5987) -- (0.9383,0.5987) -- (0.9444,0.6094) -- (0.9506,0.5987) -- (0.963,0.5987) -- (0.9568,0.588) -- (0.963,0.5774) -- (0.9753,0.5774) -- (0.9815,0.588) -- (0.9877,0.5774) -- (1,0.5774) -- (0.9938,0.5667) -- (1,0.556) -- (0.9877,0.556) -- (0.9815,0.5453) -- (0.9877,0.5346) -- (1,0.5346) -- (0.9938,0.5239) -- (1,0.5132) -- (0.9877,0.5132) -- (0.9815,0.5025) -- (0.9753,0.5132) -- (0.963,0.5132) -- (0.9568,0.5025) -- (0.963,0.4918) -- (0.9506,0.4918) -- (0.9444,0.4811) -- (0.9506,0.4704) -- (0.963,0.4704) -- (0.9568,0.4597) -- (0.963,0.4491) -- (0.9753,0.4491) -- (0.9815,0.4597) -- (0.9877,0.4491) -- (1,0.4491) -- (0.9938,0.4384) -- (1,0.4277) -- (0.9877,0.4277) -- (0.9815,0.417) -- (0.9877,0.4063) -- (1,0.4063) -- (0.9938,0.3956) -- (1,0.3849) -- (0.9877,0.3849) -- (0.9815,0.3742) -- (0.9753,0.3849) -- (0.963,0.3849) -- (0.9568,0.3742) -- (0.963,0.3635) -- (0.9506,0.3635) -- (0.9444,0.3528) -- (0.9383,0.3635) -- (0.9259,0.3635) -- (0.9321,0.3742) -- (0.9259,0.3849) -- (0.9136,0.3849) -- (0.9074,0.3742) -- (0.9012,0.3849) -- (0.8889,0.3849) -- (0.8827,0.3742) -- (0.8889,0.3635) -- (0.8765,0.3635) -- (0.8704,0.3528) -- (0.8765,0.3421) -- (0.8889,0.3421) -- (0.8827,0.3314) -- (0.8889,0.3208) -- (0.8765,0.3208) -- (0.8704,0.3101) -- (0.8642,0.3208) -- (0.8519,0.3208) -- (0.8457,0.3101) -- (0.8519,0.2994) -- (0.8395,0.2994) -- (0.8333,0.2887) -- (0.8395,0.278) -- (0.8519,0.278) -- (0.8457,0.2673) -- (0.8519,0.2566) -- (0.8642,0.2566) -- (0.8704,0.2673) -- (0.8765,0.2566) -- (0.8889,0.2566) -- (0.8827,0.2459) -- (0.8889,0.2352) -- (0.8765,0.2352) -- (0.8704,0.2245) -- (0.8765,0.2138) -- (0.8889,0.2138) -- (0.8827,0.2031) -- (0.8889,0.1925) -- (0.9012,0.1925) -- (0.9074,0.2031) -- (0.9136,0.1925) -- (0.9259,0.1925) -- (0.9321,0.2031) -- (0.9259,0.2138) -- (0.9383,0.2138) -- (0.9444,0.2245) -- (0.9506,0.2138) -- (0.963,0.2138) -- (0.9568,0.2031) -- (0.963,0.1925) -- (0.9753,0.1925) -- (0.9815,0.2031) -- (0.9877,0.1925) -- (1,0.1925) -- (0.9938,0.1818) -- (1,0.1711) -- (0.9877,0.1711) -- (0.9815,0.1604) -- (0.9877,0.1497) -- (1,0.1497) -- (0.9938,0.139) -- (1,0.1283) -- (0.9877,0.1283) -- (0.9815,0.1176) -- (0.9753,0.1283) -- (0.963,0.1283) -- (0.9568,0.1176) -- (0.963,0.1069) -- (0.9506,0.1069) -- (0.9444,0.0962) -- (0.9506,0.0855) -- (0.963,0.0855) -- (0.9568,0.0748) -- (0.963,0.0642) -- (0.9753,0.0642) -- (0.9815,0.0748) -- (0.9877,0.0642) -- (1,0.0642) -- (0.9938,0.0535) -- (1,0.0428) -- (0.9877,0.0428) -- (0.9815,0.0321) -- (0.9877,0.0214) -- (1,0.0214) -- (0.9938,0.0107) -- (1,0) -- (0.9877,0) -- (0.9815,-0.0107) -- (0.9753,0) -- (0.963,0) -- (0.9568,-0.0107) -- (0.963,-0.0214) -- (0.9506,-0.0214) -- (0.9444,-0.0321) -- (0.9383,-0.0214) -- (0.9259,-0.0214) -- (0.9321,-0.0107) -- (0.9259,0) -- (0.9136,0) -- (0.9074,-0.0107) -- (0.9012,0) -- (0.8889,0) -- (0.8827,-0.0107) -- (0.8889,-0.0214) -- (0.8765,-0.0214) -- (0.8704,-0.0321) -- (0.8765,-0.0428) -- (0.8889,-0.0428) -- (0.8827,-0.0535) -- (0.8889,-0.0642) -- (0.8765,-0.0642) -- (0.8704,-0.0748) -- (0.8642,-0.0642) -- (0.8519,-0.0642) -- (0.8457,-0.0748) -- (0.8519,-0.0855) -- (0.8395,-0.0855) -- (0.8333,-0.0962) -- (0.8272,-0.0855) -- (0.8148,-0.0855) -- (0.821,-0.0748) -- (0.8148,-0.0642) -- (0.8025,-0.0642) -- (0.7963,-0.0748) -- (0.7901,-0.0642) -- (0.7778,-0.0642) -- (0.784,-0.0535) -- (0.7778,-0.0428) -- (0.7901,-0.0428) -- (0.7963,-0.0321) -- (0.7901,-0.0214) -- (0.7778,-0.0214) -- (0.784,-0.0107) -- (0.7778,0) -- (0.7654,0) -- (0.7593,-0.0107) -- (0.7531,0) -- (0.7407,0) -- (0.7346,-0.0107) -- (0.7407,-0.0214) -- (0.7284,-0.0214) -- (0.7222,-0.0321) -- (0.716,-0.0214) -- (0.7037,-0.0214) -- (0.7099,-0.0107) -- (0.7037,0) -- (0.6914,0) -- (0.6852,-0.0107) -- (0.679,0) -- (0.6667,0) -- (0.6605,-0.0107) -- (0.6667,-0.0214) -- (0.6543,-0.0214) -- (0.6481,-0.0321) -- (0.6543,-0.0428) -- (0.6667,-0.0428) -- (0.6605,-0.0535) -- (0.6667,-0.0642) -- (0.6543,-0.0642) -- (0.6481,-0.0748) -- (0.642,-0.0642) -- (0.6296,-0.0642) -- (0.6235,-0.0748) -- (0.6296,-0.0855) -- (0.6173,-0.0855) -- (0.6111,-0.0962) -- (0.6173,-0.1069) -- (0.6296,-0.1069) -- (0.6235,-0.1176) -- (0.6296,-0.1283) -- (0.642,-0.1283) -- (0.6481,-0.1176) -- (0.6543,-0.1283) -- (0.6667,-0.1283) -- (0.6605,-0.139) -- (0.6667,-0.1497) -- (0.6543,-0.1497) -- (0.6481,-0.1604) -- (0.6543,-0.1711) -- (0.6667,-0.1711) -- (0.6605,-0.1818) -- (0.6667,-0.1925) -- (0.6543,-0.1925) -- (0.6481,-0.2031) -- (0.642,-0.1925) -- (0.6296,-0.1925) -- (0.6235,-0.2031) -- (0.6296,-0.2138) -- (0.6173,-0.2138) -- (0.6111,-0.2245) -- (0.6049,-0.2138) -- (0.5926,-0.2138) -- (0.5988,-0.2031) -- (0.5926,-0.1925) -- (0.5802,-0.1925) -- (0.5741,-0.2031) -- (0.5679,-0.1925) -- (0.5556,-0.1925) -- (0.5494,-0.2031) -- (0.5556,-0.2138) -- (0.5432,-0.2138) -- (0.537,-0.2245) -- (0.5432,-0.2352) -- (0.5556,-0.2352) -- (0.5494,-0.2459) -- (0.5556,-0.2566) -- (0.5432,-0.2566) -- (0.537,-0.2673) -- (0.5309,-0.2566) -- (0.5185,-0.2566) -- (0.5123,-0.2673) -- (0.5185,-0.278) -- (0.5062,-0.278) -- (0.5,-0.2887) -- (0.4938,-0.278) -- (0.4815,-0.278) -- (0.4877,-0.2673) -- (0.4815,-0.2566) -- (0.4691,-0.2566) -- (0.463,-0.2673) -- (0.4568,-0.2566) -- (0.4444,-0.2566) -- (0.4506,-0.2459) -- (0.4444,-0.2352) -- (0.4568,-0.2352) -- (0.463,-0.2245) -- (0.4568,-0.2138) -- (0.4444,-0.2138) -- (0.4506,-0.2031) -- (0.4444,-0.1925) -- (0.4321,-0.1925) -- (0.4259,-0.2031) -- (0.4198,-0.1925) -- (0.4074,-0.1925) -- (0.4012,-0.2031) -- (0.4074,-0.2138) -- (0.3951,-0.2138) -- (0.3889,-0.2245) -- (0.3827,-0.2138) -- (0.3704,-0.2138) -- (0.3765,-0.2031) -- (0.3704,-0.1925) -- (0.358,-0.1925) -- (0.3519,-0.2031) -- (0.3457,-0.1925) -- (0.3333,-0.1925) -- (0.3395,-0.1818) -- (0.3333,-0.1711) -- (0.3457,-0.1711) -- (0.3519,-0.1604) -- (0.3457,-0.1497) -- (0.3333,-0.1497) -- (0.3395,-0.139) -- (0.3333,-0.1283) -- (0.3457,-0.1283) -- (0.3519,-0.1176) -- (0.358,-0.1283) -- (0.3704,-0.1283) -- (0.3765,-0.1176) -- (0.3704,-0.1069) -- (0.3827,-0.1069) -- (0.3889,-0.0962) -- (0.3827,-0.0855) -- (0.3704,-0.0855) -- (0.3765,-0.0748) -- (0.3704,-0.0642) -- (0.358,-0.0642) -- (0.3519,-0.0748) -- (0.3457,-0.0642) -- (0.3333,-0.0642) -- (0.3395,-0.0535) -- (0.3333,-0.0428) -- (0.3457,-0.0428) -- (0.3519,-0.0321) -- (0.3457,-0.0214) -- (0.3333,-0.0214) -- (0.3395,-0.0107) -- (0.3333,0) -- (0.321,0) -- (0.3148,-0.0107) -- (0.3086,0) -- (0.2963,0) -- (0.2901,-0.0107) -- (0.2963,-0.0214) -- (0.284,-0.0214) -- (0.2778,-0.0321) -- (0.2716,-0.0214) -- (0.2593,-0.0214) -- (0.2654,-0.0107) -- (0.2593,0) -- (0.2469,0) -- (0.2407,-0.0107) -- (0.2346,0) -- (0.2222,0) -- (0.216,-0.0107) -- (0.2222,-0.0214) -- (0.2099,-0.0214) -- (0.2037,-0.0321) -- (0.2099,-0.0428) -- (0.2222,-0.0428) -- (0.216,-0.0535) -- (0.2222,-0.0642) -- (0.2099,-0.0642) -- (0.2037,-0.0748) -- (0.1975,-0.0642) -- (0.1852,-0.0642) -- (0.179,-0.0748) -- (0.1852,-0.0855) -- (0.1728,-0.0855) -- (0.1667,-0.0962) -- (0.1605,-0.0855) -- (0.1481,-0.0855) -- (0.1543,-0.0748) -- (0.1481,-0.0642) -- (0.1358,-0.0642) -- (0.1296,-0.0748) -- (0.1235,-0.0642) -- (0.1111,-0.0642) -- (0.1173,-0.0535) -- (0.1111,-0.0428) -- (0.1235,-0.0428) -- (0.1296,-0.0321) -- (0.1235,-0.0214) -- (0.1111,-0.0214) -- (0.1173,-0.0107) -- (0.1111,0) -- (0.0988,0) -- (0.0926,-0.0107) -- (0.0864,0) -- (0.0741,0) -- (0.0679,-0.0107) -- (0.0741,-0.0214) -- (0.0617,-0.0214) -- (0.0556,-0.0321) -- (0.0494,-0.0214) -- (0.037,-0.0214) -- (0.0432,-0.0107) -- (0.037,0) -- (0.0247,0) -- (0.0185,-0.0107) -- (0.0123,0) -- cycle;

          \end{scope}
          \draw[sectioncolor!45, dashed, thin] (0,0) -- (0.5,0.8660254) -- (1,0) -- cycle;
        \end{tikzpicture}}
    \end{column}
    \begin{column}{0.52\textwidth}
      \small
      Quatrième étape de la construction, superposée au triangle de
      départ (en pointillé).
      \vspace{6pt}

      À la limite : un contour de longueur infinie, qui délimite une
      région d'aire finie.
      \vspace{6pt}

      Le mot fractale : Benoît Mandelbrot, 1975, du latin fractus,
      « brisé ». Des irrégularités à toutes les échelles ;
      l'agrandissement ne les rend jamais lisses.
    \end{column}
  \end{columns}
\end{frame}

\begin{frame}{Des monstres devenus modèles}
  \small
  \begin{itemize}
    \item Plusieurs de ces objets sont connus depuis la fin du
      \siecle{19} : des contre-exemples (courbes continues sans
      tangente, ensembles troués à l'infini), qu'on appelait des
      « monstres ».
    \item L'ordinateur permet de les représenter, et de voir qu'ils
      décrivent bien des formes naturelles : côtes, nuages, végétaux,
      réseaux sanguins.
  \end{itemize}
  \vspace{6pt}
  Le chemin du chapitre :
  \begin{enumerate}
    \item deux exemples : la courbe de von Koch, le triangle de Sierpiński ;
    \item autosimilarité et dimension de similitude ;
    \item systèmes itératifs de fonctions (SIF), justifiés par Banach ;
    \item deux algorithmes de dessin ;
    \item des monstres de Weierstrass aux modèles de Mandelbrot.
  \end{enumerate}
\end{frame}
